% ============================================================================
%  CAPÍTULO I: ENTORNO DE PRUEBAS Y DEFINICIÓN DEL OBJETIVO
% ============================================================================

\section{Introducción al Entorno de Evaluación}

La evaluación de seguridad y el análisis de vulnerabilidades constituyen una disciplina fundamental dentro de la gestión de riesgos tecnológicos modernos \citep{nist_sp800_115}. El presente informe documenta de manera exhaustiva el proceso de escaneo, identificación, categorización y mitigación de fallos de seguridad sobre la infraestructura del servidor corporativo de la organización financiera \textit{Banco Sudoers}.

El ejercicio se realizó bajo la modalidad de auditoría de caja gris (\textit{Gray-Box Assessment}), donde el equipo auditor posee conocimiento preliminar del direccionamiento IP y la tipología de servicios desplegados, pero simula las técnicas de descubrimiento y vectores de ataque que ejecutaría un actor malicioso externo o interno con acceso al segmento de red local.

\section{Definición del Equipo Objetivo (Target)}

El equipo objetivo auditado corresponde al servidor central de operaciones e infraestructura (\texttt{dc1.sudoers.lan}), el cual consolida los servicios críticos de negocio, autenticación, almacenamiento, correo electrónico y mecanismos defensivos de detección de intrusos.

\subsection{Ficha Técnica del Sistema Objetivo}

En la \textbf{TABLA \ref{tab:ficha-target}} se especifican los parámetros técnicos, hardware y de red que caracterizan al host auditado durante la ventana de evaluación de seguridad.

\begin{table}[H]
\centering
\caption{Ficha Técnica del Servidor Objetivo (Target)}
\label{tab:ficha-target}
\contenidoConFuente{%
  \begin{tablaAPA}
  \begin{tabular}{L{4.2cm}L{10.8cm}}
  \toprule
  \textbf{Parámetro} & \textbf{Valor / Descripción Técnica} \\
  \midrule
  \textbf{Nombre de Host (FQDN)} & \texttt{dc1.sudoers.lan} (Alias de red: \texttt{target-host}) \\
  \textbf{Dirección IPv4 del Target} & \texttt{192.168.100.85} (Asignada en interfaz LAN cableada/Wi-Fi) \\
  \textbf{Máscara de Red / Prefijo} & \texttt{255.255.255.0} (\texttt{/24}) \\
  \textbf{Puerta de Enlace (Gateway)} & \texttt{192.168.100.1} \\
  \textbf{Sistema Operativo Base} & GNU/Linux Debian 12 / Kali Linux Rolling (Kernel 6.x x86\_64) \\
  \textbf{Modelo Operativo} & Arquitectura híbrida: Servicios de aplicación en contenedores Docker y motores de defensa nativos en espacio de kernel \\
  \textbf{Rol Corporativo} & Servidor Central de Banca Digital, Servicios Intranet y Pasarela Perimetral \\
  \textbf{Dirección MAC del Target} & \texttt{08:00:27:C3:FB:76} (Oracle VirtualBox virtual NIC) \\
  \textbf{Estación de Auditoría} & Estación Kali Linux dedicada (\texttt{192.168.100.50}) \\
  \bottomrule
  \end{tabular}
  \end{tablaAPA}%
}{Elaboración propia con base en el levantamiento de infraestructura.}
\end{table}

\subsection{Matriz de Servicios y Puertos Expuestos}

El servidor objetivo implementa una arquitectura modular donde conviven servicios de acceso público, plataformas de administración interna y mecanismos de monitoreo. La \textbf{TABLA \ref{tab:matriz-servicios}} detalla la distribución de puertos que fueron sometidos al proceso de escaneo y análisis de vulnerabilidades.

\begin{table}[H]
\centering
\caption{Matriz de Servicios y Puertos en Evaluación}
\label{tab:matriz-servicios}
\contenidoConFuente{%
  \begin{tablaAPA}
  \begin{tabular}{C{1.6cm}C{1.8cm}L{3.8cm}L{3.6cm}L{3.8cm}}
  \toprule
  \textbf{Puerto} & \textbf{Protocolo} & \textbf{Servicio / Software} & \textbf{Entorno de Ejecución} & \textbf{Propósito en la Organización} \\
  \midrule
  80 & TCP & Nginx Reverse Proxy & Contenedor Docker & Portal Web Financiero y Redirección \\
  443 & TCP & Nginx TLS Server & Contenedor Docker & Acceso Seguro HTTPS Bancario \\
  8080 & TCP & Roundcube Webmail & Contenedor Docker & Plataforma de Correo Corporativo \\
  8082 & TCP & TinyFileManager & Contenedor Docker & Gestor Web de Archivos y Documentos \\
  3000 & TCP & Grafana Analytics & Contenedor Docker & Tableros de Telemetría y Métricas \\
  5432 & TCP & PostgreSQL Engine & Contenedor Docker & Base de Datos Relacional de Clientes \\
  5636 & TCP & EveBox GUI & Contenedor Docker & Consola de Alertas IDS/IPS en Tiempo Real \\
  9392 & TCP & Greenbone GSA & Contenedor Docker & Consola de Auditoría y Escaneo OpenVAS \\
  22 & TCP & OpenSSH Daemon & Host Nativo Linux & Administración Remota Segura por Consola \\
  \bottomrule
  \end{tabular}
  \end{tablaAPA}%
}{Elaboración propia con base en la configuración de \texttt{docker-compose.yml} y el firewall del host.}
\end{table}

\newpage
\section{Topología de Red del Escenario de Pruebas}

Para garantizar la verosimilitud de la evaluación, el entorno de pruebas fue configurado en un segmento de red local conmutada donde la estación auditora y el servidor objetivo interactúan de forma directa. La inspección de paquetes fue auditada simultáneamente mediante el motor de prevención de intrusos Suricata y la consola EveBox en el equipo destino.

\begin{figure}[H]
\centering
\caption{Diagrama de Topología del Entorno de Auditoría y Escaneo}
\label{fig:topologia-pruebas}
\resizebox{0.96\linewidth}{!}{%
\begin{tikzpicture}[
  box/.style={rectangle, draw=black!80, thick, fill=gray!8, align=center, rounded corners=4pt, inner sep=8pt, minimum height=2.6cm, minimum width=3.6cm, font=\small\sffamily},
  targetbox/.style={rectangle, draw=red!75!black, very thick, fill=red!4, align=center, rounded corners=4pt, inner sep=8pt, minimum height=2.6cm, minimum width=5.2cm, font=\small\sffamily},
  scannerbox/.style={rectangle, draw=blue!75!black, very thick, fill=blue!4, align=center, rounded corners=4pt, inner sep=8pt, minimum height=2.6cm, minimum width=4.4cm, font=\small\sffamily},
  line/.style={draw, thick, -{Latex[length=2.8mm]}},
  dashline/.style={draw, thick, dashed, red!75!black, -{Latex[length=2.8mm]}}
]
  % Nodos principales
  \node[scannerbox] (auditor) {
    \textbf{Estación Auditora}\\[2pt]
    \footnotesize (Kali Linux Rolling)\\[4pt]
    \footnotesize IP: \texttt{192.168.100.50}\\[3pt]
    \scriptsize Nmap, OWASP ZAP, Wireshark
  };

  \node[box, right=2.5cm of auditor] (switch) {
    \textbf{Switch / Router LAN}\\[4pt]
    \footnotesize Red: \texttt{192.168.100.0/24}\\[3pt]
    \scriptsize Gateway: \texttt{192.168.100.1}
  };

  \node[targetbox, right=2.5cm of switch] (target) {
    \textbf{Servidor Objetivo (Target)}\\[2pt]
    \footnotesize \texttt{dc1.sudoers.lan}\\[2pt]
    \footnotesize (\texttt{192.168.100.85})\\[4pt]
    \scriptsize Nginx, PostgreSQL, Webmail, Filemanager\\[2pt]
    \scriptsize \textbf{Defensas:} Nftables + Suricata IDS/IPS
  };

  % Líneas de conexión con etiquetas multilínea limpias
  \path[line] ([yshift=0.35cm]auditor.east) -- node[above=2pt, font=\scriptsize\sffamily, align=center] {Tráfico de\\Evaluación} ([yshift=0.35cm]switch.west);
  \path[line] ([yshift=0.35cm]switch.east) -- node[above=2pt, font=\scriptsize\sffamily, align=center] {Sondas de\\Escaneo} ([yshift=0.35cm]target.west);
  \path[dashline] ([yshift=-0.55cm]target.west) -- node[below=2pt, font=\scriptsize\sffamily, align=center] {Alertas IDS /\\Drop IPS} ([yshift=-0.55cm]switch.east);
\end{tikzpicture}%
}
\fuenteAPA{Elaboración propia.}
\end{figure}

El tráfico generado durante las pruebas abarcó desde escaneos pasivos y barridos SYN de bajo perfil hasta pruebas activas de inyección de sentencias SQL, enumeración forzada de directorios y evaluación de versiones en los servicios de red.
